\section*{Chapter \ref{ch:b3:ode} --- Ordinary Differential
Equations}

\begin{solution}{exo:b3:ode:1}
(a) Separating variables: $x(t) = \frac1{1 - t}$ on
$\intoo{-\infty}1$: blow-up at $T_+ = 1$, with $x(t) \to
+\infty$ --- exactly \cref{thm:b3:ode:maximal}(2).
(b) $x(t) = \tan t$ on $\intoo{-\pi/2}{\pi/2}$: blow-up at
both ends.
(c) $x(t) = \frac{1}{1 + \eu^{-t}}$, global: the solution
stays in $\intoo01$, a bounded set, so the graph cannot
escape every compact of $\R\times\R$ in finite time --- $T_\pm
= \pm\infty$. Note (a), (b) do not contradict
\cref{cor:b3:ode:globallinear}: $x^2$ and $1 + x^2$ have
superlinear growth.
\end{solution}

\begin{solution}{exo:b3:ode:2}
(a) This is \cref{cor:b3:ode:uniqueness}'s estimate with $t_0
= 0$. Sharpness: for $x' = Lx$ (globally $L$-Lipschitz), two
solutions differ by exactly $(x_0 - y_0)\eu^{Lt}$.
(b) The estimate reads: the time-$t$ flow map is
$\eu^{L\abs t}$-Lipschitz in the initial condition ---
continuity, uniformly for $t$ in compacts; on bounded sets of
non-globally-Lipschitz $f$, run the same on a compact tube
around the trajectories with the local constant.
\end{solution}

\begin{solution}{exo:b3:ode:3}
(a) $\abs{\sin(tx)} \leq 1$: bounded, i.e.\ linear growth with
$\alpha = 0$, $\beta = 1$: \cref{cor:b3:ode:globallinear} on
$U = \R\times\R$.
(b) $\abs{tx/(1 + x^2)} \leq \abs t\cdot\frac12$: again
sublinear (in fact bounded on time-compacts): global.
(c) $X = (x, x')$: $X' = \bigl(\begin{smallmatrix}0 & 1\\ -q(t)
& 0\end{smallmatrix}\bigr)X$: linear with continuous
coefficients: global (\cref{thm:b3:ode:linearstructure}'s
setting).
(d) Global; Gr\"onwall as in \cref{cor:b3:ode:globallinear}:
$\norm{x(t)} \leq \norm{x(t_0)}\,\eu^{M\abs{t - t_0}}$ with
$M = \sup\vertiii{A}$.
\end{solution}

\begin{solution}{exo:b3:ode:4}
(a) $A^2 = -I$ for the first: $\eu^{tA} = \cos t\,I + \sin
t\,A = \bigl(\begin{smallmatrix}\cos t & -\sin t\\ \sin t &
\cos t\end{smallmatrix}\bigr)$. Jordan block: $\lambda I$ and
$N$ commute: $\eu^{tA} =
\eu^{\lambda t}\bigl(\begin{smallmatrix}1 & t\\ 0 &
1\end{smallmatrix}\bigr)$. Third: $A^2 = I$: $\eu^{tA} =
\cosh t\,I + \sinh t\,A$.

(b) System $X' = \bigl(\begin{smallmatrix}0&1\\-1&0
\end{smallmatrix}\bigr)X + \bigl(\begin{smallmatrix}0\\
\cos\omega t\end{smallmatrix}\bigr)$; Duhamel with the
rotation resolvent gives the particular solutions: for
$\omega \neq 1$, $x_p(t) = \frac{\cos(\omega t)}{1 -
\omega^2}$ (verify directly); for $\omega = 1$ the integral
$\int_0^t\sin(t - s)\cos s\,\dd s = \frac t2\sin t$ produces
the \emph{secular} growth $x_p = \frac t2\sin t$: resonance
--- the forcing pumps energy at the natural frequency and the
amplitude grows linearly.
\end{solution}

\begin{solution}{exo:b3:ode:5}
(a) $w' = x_1x_2'' - x_1''x_2 = x_1(-px_2' - qx_2) - (-px_1'
- qx_1)x_2 = -p\,w$: $w(t) = w(t_0)\eu^{-\int_{t_0}^tp}$,
never zero or identically zero. If $w \neq 0$, the vectors
$(x_i, x_i')(t_0)$ are independent in $\R^2$, and since the
solution space has dimension $2$
(\cref{thm:b3:ode:linearstructure} for the system), $(x_1,
x_2)$ is a basis.

(b) With $u = \int\frac{\eu^{-\int p}}{x_1^2}$: $x_2 = x_1u$,
$x_2' = x_1'u + \frac{\eu^{-\int p}}{x_1}$, and
\[
x_2'' + px_2' + qx_2
= u\,(x_1'' + px_1' + qx_1) +
\Bigl(-\,p\frac{\eu^{-\int p}}{x_1} +
p\frac{\eu^{-\int p}}{x_1}\Bigr) = 0
\]
(the cross terms cancel exactly; expand carefully). For
$t^2x'' - 2x = 0$, i.e.\ $x'' - \frac2{t^2}x = 0$ ($p = 0$)
with $x_1 = t^2$: $u = \int t^{-4} = -\frac1{3t^3}$, so $x_2
= -\frac1{3t}$: general solution $at^2 + \frac bt$.
\end{solution}

\begin{solution}{exo:b3:ode:6}
Equilibria $0, 1$; $f'(x) = 1 - 2x$: $f'(0) = 1 > 0$
(unstable --- nearby solutions move away, as the explicit
form shows), $f'(1) = -1 < 0$: asymptotically stable
(\cref{thm:b3:ode:linearization} in dimension $1$). For
$x(0) \in \intoo01$: $f > 0$ there, so as long as the
solution stays in $\intoo01$ it increases; it can never reach
$0$ or $1$ (uniqueness: those are trajectories), so it stays,
is bounded --- hence global --- and increases to a limit $L
\in \intoc{x(0)}1$. If $f(L) \neq 0$, then $x' \geq c > 0$
near the limit, forcing $x$ past $L$: so $f(L) = 0$, $L = 1$;
symmetrically $x \to 0$ at $-\infty$. Explicitly $x(t) =
\frac1{1 + C\eu^{-t}}$ confirms everything. Generally: if a
scalar autonomous solution had $x'(t_0) = 0$, then $x(t_0)$
is an equilibrium and uniqueness makes $x$ constant;
otherwise $f(x(t))$ keeps a fixed sign (it never vanishes,
and $t \mapsto f(x(t))$ is continuous): $x$ is strictly
monotone --- so a nonconstant periodic solution is
impossible.
\end{solution}

\begin{solution}{exo:b3:ode:7}
(a) $\frac{\dd}{\dd t}H(x,y) = H_xx' + H_yy' = H_xH_y +
H_y(-H_x) = 0$.
(b) The level sets of $H = \frac{y^2}2 + \frac{x^4}4$ are
compact ($H$ coercive), so solutions are trapped in compacts:
global and bounded (\cref{thm:b3:ode:maximal}). Stability of
$(0,0)$: $V = H$ is positive definite ($H = 0$ only at the
origin) with $\dot V = 0$: \cref{thm:b3:ode:lyapunov}. The
linearization $x' = y$, $y' = 0$ has the non-diagonalizable
nilpotent matrix with eigenvalue $0$:
\cref{thm:b3:ode:linearization} is silent (its hypothesis
$\operatorname{Re}\lambda < 0$ fails), and indeed the
linearized system is unstable ($y_0 \ne 0$ drifts) while the
nonlinear one is stable: linearization at a non-hyperbolic
equilibrium proves nothing.
\end{solution}

\begin{solution}{exo:b3:ode:8}
(a) $\dot V = y\,y' + \sin x\cdot x' = y(-\sin x - cy) +
y\sin x = -cy^2 \leq 0$, and $V = \frac{y^2}2 + (1 - \cos x)$
is positive definite on $\{\abs x < 2\pi\}$ around the
origin: stable (\cref{thm:b3:ode:lyapunov}).
(b) The linearized matrix at $(0,0)$ is
$\bigl(\begin{smallmatrix}0 & 1\\ -1 &
-c\end{smallmatrix}\bigr)$, with characteristic polynomial
$\lambda^2 + c\lambda + 1$: roots $\frac{-c \pm \sqrt{c^2 -
4}}2$ --- both real negative if $c \geq 2$, complex with real
part $-\frac c2 < 0$ if $0 < c < 2$. In all cases
$\operatorname{Re}\lambda < 0$:
\cref{thm:b3:ode:linearization} gives asymptotic stability
(despite the degenerate $\dot V$).
(c) At $(\pi, 0)$: $\sin(\pi + u) = -\sin u$, linearization
$\bigl(\begin{smallmatrix}0&1\\1&-c\end{smallmatrix}\bigr)$,
characteristic $\lambda^2 + c\lambda - 1$: roots of opposite
signs ($\lambda_+\lambda_- = -1$). Along the unstable
eigenvector, the linear system has the explicitly escaping
solution $\eu^{\lambda_+t}v_+$ with $\lambda_+ > 0$: the
inverted pendulum is unstable for every damping.
\end{solution}

\begin{solution}{exo:b3:ode:9}
For $\alpha \in \intoo01$: besides $x \equiv 0$, each
\[
x_c(t) = \begin{cases}0 & t \leq c,\\
\bigl((1-\alpha)(t - c)\bigr)^{1/(1-\alpha)} & t \geq c,
\end{cases}
\]
is $\mathcal C^1$ and solves the equation (the exponent
$\frac1{1-\alpha} > 1$ makes the derivative vanish at $c$): a
continuum of solutions through $(0,0)$. For $\alpha = 1$: $x
\mapsto \abs x$ is globally $1$-Lipschitz
($\abs{\abs a - \abs b} \leq \abs{a - b}$), so
\cref{thm:b3:ode:cauchylipschitz} applies and the only
solution through $0$ is $x \equiv 0$. The frontier is exactly
the local Lipschitz condition at $0$: $\abs x^\alpha$ has
unbounded difference quotients there for $\alpha < 1$.
\end{solution}

\begin{solution}{exo:b3:ode:10}
(a) Suppose $\norm{x(t_2)} > R$ for some $t_2 > 0$ with
$\norm{x(0)} \leq R$, and let $t_1 = \sup\{t \leq t_2 :
\norm{x(t)} \leq R\}$: then $\norm{x(t_1)} = R$ and
$\norm{x(t)} > R$ on $\intoc{t_1}{t_2}$. On that interval
$g(t) = \norm{x(t)}^2$ has $g'(t) = 2\langle x(t),
F(x(t))\rangle \leq 0$ (the hypothesis applies: $\norm{x(t)}
\geq R$), so $g(t_2) \leq g(t_1) = R^2$: contradiction. The
ball is positively invariant.
(b) A solution trapped in the compact ball cannot have
$T_+ < \infty$ (\cref{thm:b3:ode:maximal}(2)): global
forward. Gradient system: $\frac{\dd}{\dd t}G(x(t)) =
\langle\nabla G, -\nabla G\rangle = -\norm{\nabla G(x(t))}^2
\leq 0$: $G$ decreases, so the solution stays in $\{G \leq
G(x_0)\}$, which is bounded (coercivity: outside a large
ball, $G > G(x_0)$) and closed: compact. Escape is
impossible: every solution of a coercive gradient system is
global forward, sliding downhill forever.
\end{solution}

\begin{solution}{pb:b3:ode:1}
\textbf{1.} $\dot E = yy' + \sin x\cdot x' = -y\sin x +
y\sin x = 0$: energy is a first integral. On a maximal
solution, $y^2 = 2(E_0 + \cos x) \leq 2(E_0 + 1)$: $y$ is
bounded; then $\abs{x(t)} \leq \abs{x(0)} +
t\sup\abs y$ grows at most linearly: on any finite time
interval the trajectory stays in a compact of $\R^2$, so
\cref{thm:b3:ode:maximal}(2) forces $T_\pm = \pm\infty$.
Equilibria $(k\pi, 0)$; linearization
$\bigl(\begin{smallmatrix}0&1\\
\mp1&0\end{smallmatrix}\bigr)$ with $-\cos(k\pi) = \mp1$:
eigenvalues $\pm\iu$ for even $k$ (center type, inconclusive
by itself) and $\pm1$ for odd $k$ (saddle).

\textbf{2.} $E$ constant along solutions confines each
trajectory to a level set $\{y^2 = 2(E_0 + \cos x)\}$. For
$E_0 = -1$: only the points $(2k\pi, 0)$. For $-1 < E_0 <
1$: writing $E_0 = -\cos a$, the set is a disjoint union of
closed curves $y = \pm\sqrt{2(\cos x - \cos a)}$ over $x \in
[2k\pi - a, 2k\pi + a]$, one around each stable equilibrium.
For $E_0 = 1$: the curves $y = \pm2\cos\frac x2$ joining
consecutive saddles --- the separatrix --- together with the
saddles themselves. For $E_0 > 1$: two graphs $y =
\pm\sqrt{2(E_0 + \cos x)}$, defined for all $x$, never
touching $y = 0$.

\textbf{3.} $V = E + 1 = \frac{y^2}2 + (1 - \cos x)$
vanishes at $(0,0)$, is positive on a punctured neighborhood
($\abs x < 2\pi$), and $\dot V = 0 \leq 0$:
\cref{thm:b3:ode:lyapunov} gives stability. Not asymptotic:
the solution through $(a, 0)$ ($0 < a$ small) remains on the
level curve $E = -\cos a$, whose distance to the origin is
positive (the curve meets the $x$-axis only at $\pm a$):
$\varphi_t(a, 0) \not\to (0,0)$.

\textbf{4.} On the level curve $C_a$: no equilibria ($y = 0$
forces $x = \pm a$ with $\sin(\pm a) \neq 0$ for $0 < a <
\pi$), so the speed $\norm{(y, -\sin x)}$ has a positive
minimum $m$ on the compact $C_a$. Follow the solution from
$(a, 0)$: in the lower half plane $x' = y < 0$, so $x$
decreases from $a$ to $-a$ in finite time (the quarter/half
period integrals converge: question 5's analysis), reaching
$(-a, 0)$; by the symmetry $(x, y) \mapsto (x, -y)$, $t
\mapsto -t$ of the equation, the upper half is traversed
back in the same time $\frac T2$: the solution returns to
$(a, 0)$ at time $T$. Uniqueness
(\cref{cor:b3:ode:uniqueness}) then propagates: $x(t + T) =
x(t)$ for all $t$: periodic.

\textbf{5.} On the branch where $y > 0$: $y = \sqrt{2(\cos x
- \cos a)}$ and $\dd t = \frac{\dd x}{y}$; integrating $x$
from $-a$ to $a$ gives the half period, and the symmetry $x
\mapsto -x$ halves the integral again:
\[
T(a) = 4\int_0^a\frac{\dd x}{\sqrt{2(\cos x - \cos a)}} .
\]
Convergence at $x = a^-$: $\cos x - \cos a = \sin(a)(a - x) +
O((a-x)^2)$ with $\sin a > 0$: the integrand behaves like
$\bigl(2\sin a\,(a - x)\bigr)^{-1/2}$, integrable.

\textbf{6.} With $\sin\frac x2 = k\sin\varphi$, $k =
\sin\frac a2$: $\cos x - \cos a = 2(k^2 - \sin^2\frac x2)
= 2k^2\cos^2\varphi$ and $\frac12\cos\frac x2\,\dd x =
k\cos\varphi\,\dd\varphi$, so
\[
T(a) = 4\int_0^{\pi/2}
\frac{\dd\varphi}{\sqrt{1 - k^2\sin^2\varphi}} .
\]
As $a \to 0^+$, $k \to 0$: for $k \leq k_0 < 1$ the
integrand is dominated by $(1 - k_0^2\sin^2\varphi)^{-1/2}$,
continuous on $\intcc0{\pi/2}$: DCT gives $T \to
4\cdot\frac\pi2 = 2\pi$. Expanding $(1 - u)^{-1/2} = 1 +
\frac u2 + \frac{3u^2}8 + \cdots$ with $u =
k^2\sin^2\varphi$ (normal convergence for $k < 1$) and
$\int_0^{\pi/2}\sin^2 = \frac\pi4$:
\[
T(a) = 2\pi\Bigl(1 + \frac{k^2}{4} + O(k^4)\Bigr) :
\]
isochronism holds only to first order; the period grows with
amplitude.

\textbf{7.} As $k \uparrow 1$ the integrands increase to $(1
- \sin^2\varphi)^{-1/2} = \frac1{\cos\varphi}$, whose
integral diverges: by monotone convergence, $T(a) \to
+\infty$ as $a \to \pi^-$.

\textbf{8.} On the branch $y = 2\cos\frac x2$ ($\abs x <
\pi$): $\frac{\dd x}{2\cos(x/2)} = \dd t$; with $u = \frac
x2$, $\int\frac{\dd u}{\cos u} =
\ln\tan\bigl(\frac u2 + \frac\pi4\bigr)$, so $t =
\ln\tan\bigl(\frac x4 + \frac\pi4\bigr)$, i.e.
\[
x(t) = 4\arctan(\eu^t) - \pi,
\qquad y(t) = x'(t) = \frac{4\eu^t}{1 + \eu^{2t}} =
\frac{2}{\cosh t} .
\]
Verification via the energy: with $\theta = \arctan\eu^t$,
$\sin2\theta = \frac1{\cosh t}$, so $\cos x = -\cos4\theta =
-1 + \frac{2}{\cosh^2t}$ and $E = \frac{y^2}2 - \cos x =
\frac2{\cosh^2t} + 1 - \frac2{\cosh^2t} = 1$: the trajectory
lies on the separatrix, and differentiating $y^2 = 2(1 +
\cos x)$ where $y > 0$ reproduces $y' = -\sin x$. Limits: $x
\to \pm\pi$ and $y \to 0$ as $t \to \pm\infty$.

\textbf{9.} The orbit tends to the saddle $(\pi, 0)$ forward
and $(-\pi, 0)$ backward but never arrives: if it reached
$(\pi, 0)$ at a finite time $t^*$, two distinct maximal
solutions --- the separatrix solution and the constant
solution at the saddle --- would pass through the same point
$(t^*, (\pi, 0))$, contradicting
\cref{cor:b3:ode:uniqueness}. Saddles are approached only
asymptotically.

\textbf{10.} For $E_0 > 1$: $y^2 = 2(E_0 + \cos x) \geq
2(E_0 - 1) > 0$: $y$ keeps its sign, and $\abs{x'} = \abs y
\geq \sqrt{2(E_0 - 1)}$: $x$ is strictly monotone, global
(question 1), with $x(t) \to \pm\infty$. Since $y(t) =
\pm\sqrt{2(E_0 + \cos x(t))}$ and $\cos$ is $2\pi$-periodic,
$y$ returns to its value each time $x$ advances by $2\pi$;
the time needed is
\[
\tau(E_0) = \int_{x_0}^{x_0 + 2\pi}\frac{\dd
x}{\sqrt{2(E_0 + \cos x)}} =
\int_{-\pi}^{\pi}\frac{\dd x}{\sqrt{2(E_0 + \cos x)}}
\]
(substitution; periodicity): $y$ is $\tau$-periodic ---
the pendulum whirls with asymptotically constant rotation
rate $2\pi/\tau \approx \sqrt{2E_0}$ for large energies.

\textbf{11.} The method, in order: \emph{energy} ($\dot E =
0$) reduces the two-dimensional flow to one-dimensional
level curves; \emph{boundedness} of $y$ on each level plus
escape-from-compacts gives global existence;
\emph{compactness} of the closed levels gives speed bounds
and hence periodicity; \emph{uniqueness} converts the
first-return into exact periodicity, forbids finite-time
arrival at saddles, and separates the orbit types;
\emph{linearization and Lyapunov} classify the equilibria;
the \emph{period integral} is analyzed with the convergence
theorems of \cref{ch:b3:lebesgue}. At no point did we
possess --- or need --- a closed-form general solution: the
qualitative theory extracted every feature of the motion
from the equation itself.
\end{solution}
